\section{Energy integration}


%Nasibeh use cetow for heat puming solutions and st fons for large scale integration and combined heat and power production. We have to use neutralised values (i.e. we start with 100 unit and end up with 46.


%Converting energy resources into useful energy, exergy efficiency, heat pumping and combined heat and power => examples from the papers of Nasibeh

\subsection{Multi-level energy requirement definition}
\label{subsec:Multi-layer}
The necessary data for the energy requirement analysis are the heating or cooling enthalpy demand $(\dot{Q})$, and the inlet and outlet temperature $ (T) $ of the PUOs. In energy analysis, these requirements are defined as the Temperature-Enthalpy (T-H) profiles for each of those process units. Based on the energy requirement definition, the data extraction can be categorized into five levels, with the following description. (see Table.\ref{tab2:table0})

\begin{itemize}
	\item \textbf{Black-box analysis:}
        It considers the global view of the site where the result comes from the analysis of the current energy bill of the plant. In this level, the process-utility heat transfer units are considered from their utility interface so that both the temperature and the heat load of the energy requirements is represented by their corresponding utility data.
    \item \textbf{Grey-box analysis:}
        This analysis updates the temperature level of the energy requirements to its corresponding process data by considering the process interface of the process-utility heat transfer unit. At this level, all the process-process (P-P) heat exchangers are accepted as their existing configuration. This is the typical level of data extraction in Total Site Analysis.
    \item \textbf{White-box analysis:}
         In this third level, in addition to the previous features, the process-process heat transfer units are also considered. 
    \item \textbf{Simple-model analysis:}
    	In this level, the analysis is additionally performed for non-isothermal mixers and reactors (including the streams that leave or enter into the system with the potential of heat recovery) based on the simple model. 
    \item \textbf{Detailed-model analysis:}
    	In the last level of analysis, all necessary detailed process information is extracted to define the T-H profiles of requirements from their corresponding process data using advanced simulation tool. The non-linear T-H profiles of streams with variable specific heat ($C_{p}$) are recognizable at this level.
\end{itemize}



\definecolor{bell1}{RGB}{77,77,77}
\definecolor{bell2}{RGB}{153,153,153}
\definecolor{bell3}{RGB}{230,230,230}

\begin{table*}[h] 
\caption{Multi-Level data extraction \cite{pouransari2014}}
\label{tab2:table0} \vspace*{-2mm}
\begin{center}
\resizebox{\textwidth}{!}
{
$\rotatebox[origin=t]{270}{\textbf{Top-Down analysis}}%
%LEFTARROW--------------------------------------------------
\hspace{0.3cm}\raisebox{-2.5cm}{\begin{tikzpicture}[>=latex]
\draw[<-,black,line width=0.07cm] (0,-4.95) -- (0,0);
\end{tikzpicture}}
%LEFTARROW--------------------------------------------------
%\left\downarrow% Use `\left.` if don't want arrow on this side.
\resizebox{14cm}{!} {
\hspace{0.3cm}\begin{tabular} { >{\centering} m{1.3cm} >{\centering}  m{1.7cm} >{\centering} m{1.7cm} >{\centering} m{1.9cm} >{\centering} m{3.8cm} >{\centering} m{3.2cm} >{\centering\arraybackslash} m{4.2cm} }
\noalign{\hrule height 0.08cm}
 \vspace*{1mm} {\textbf{Level}} &\vspace*{1mm} \textbf{$T$ } & \vspace*{1mm}\textbf{$\dot{Q}$} & \vspace*{1mm}\textbf{P-P HEXs } & \vspace*{1mm}\textbf{Non-Isothermal mixers \& reactors } & \vspace*{1mm}\textbf{Analysis Type } & \vspace*{1mm}\textbf{Energy requirement representation } \\ [0.1cm]
\hline \vspace*{2mm}
1 & \vspace*{1mm} Utility & \vspace*{1mm} Utility& \vspace*{1mm} $\boldsymbol{\timesbig}$  & \vspace*{1mm} $\boldsymbol{\timesbig}$ & \vspace*{1 mm}  \cellcolor{black} \textbf{\textcolor{white}{Black-Box}}  & \vspace*{5 mm}  Utility \\  [0.3 cm]
& & & & & \\  [-0.35cm]
2  &  Process & Utility& $\boldsymbol{\timesbig }$  & $\boldsymbol{\timesbig}$ & \cellcolor{bell1} \textbf{\textcolor{white}{ Grey-Box}}  & Technology, Thermodynamic   \\ [0.1cm]
& & & & & \\  [-0.35cm]
3  & Process & Utility & $\boldsymbol{\checkmarkbig}$ & $\boldsymbol{\timesbig}$ & \cellcolor{bell2} White-Box (Deatiled-Pinch)  & Technology, Thermodynamic    \\ [0.1cm]
& & & & & \\  [-0.35cm]
4  & Process & Utility & $\boldsymbol{\checkmarkbig}$ & $\boldsymbol{\checkmarkbig}$ & \cellcolor{bell3} Simple-Model  & Technology, Thermodynamic   \\ [0.1cm]
& & & & & \\  [-0.35cm]
5  & Process& Process & $\boldsymbol{\checkmarkbig}$ & $\boldsymbol{\checkmarkbig}$ &  Detailed-Model & Technology, Thermodynamic   \\ [0.1cm]
\noalign{\hrule height 0.08cm}
\end{tabular}}
%\right\Uparrow %  Use `\right.` if don't want arrow on this side.
%RIGHTARROW--------------------------------------------------
\hspace*{0.3cm}\raisebox{-2.5cm}{\begin{tikzpicture}[>=latex]
\draw[<-,black,line width=0.07cm] (0,0) -- (0,-4.95);
\end{tikzpicture}}
%RIGHTARROW--------------------------------------------------
\hspace{0.3cm}\rotatebox[origin=b]{90}{\textbf{Bottom-Up analysis}}$
}
\end{center}
\end{table*}

\subsection{Multiple representation of the same energy requirement}
\label{subsec2:Multi-rep}

In the process analysis methodology of \citet{brown2005dual}, \citet{muller2007energy} and then in \citet{Perin-Levasseur2009}, the energy requirement of each process unit is classified into three representations: \textbf{\textit{utility}}, \textbf{\textit{technology}} and \textbf{\textit{thermodynamic}} representation. The utility requirement would take the energy requirement of a unit equivalent to the production of the distributed energy used in the unit. Thus, the sum of the utility requirements defines the steam consumption of the factory. In this case neither the process nor the utility supply can be modified. Therefore, including one stream with utility representation means there is no need to modify the current heat exchangers serve for the same heat demand related to that stream. The technology requirement can be defined by the heat exchange interface between the technology that realizes the unit operation and the distributed energy. It represents the T-H profile required by the process streams in existing heat exchangers. With technology representation only heat supply through the heat transfer unit can be modified, but not the process unit itself. The thermodynamic requirement corresponds to the heat transfer required to perform a unit operation. The definition results from the analysis of the thermo-chemical and physical transformations to be realized in the PUO. The thermodynamic requirement describes the necessary thermodynamic transformations in the operation and defines the lowest temperature profile for the cold stream and the highest temperature profile for the hot stream. This representation is chosen when the process and the heat supply can be modified. With the assumption of no heat losses, the representations are characterized by the same amount of the heat, but with different temperature profiles. They have the same energy content but a  different exergy content.  In the context of process integration, the exergy concept is combined with the pinch analysis to reduce the energy requirement of the process \cite{LENI-ARTICLE-1996-002}. Here, the exergy concept has been applied to differentiate the heat transfer units that serve for the same energy requirement. Now we make the link between the two concepts: each of the five mentioned data extraction levels can be explained by one or two representation types, which is summarized in the right side of the Table. \ref{tab2:table0}. The black-box level corresponds to the utility representation. The grey-box to the simple-model levels correspond to the technology or the thermodynamic representation and the detailed-model level defines the thermodynamic requirement of the PUOs. 

\begin{figure*}[!ht]
 \centering
  \begin{tabular}{c}
 	\subfloat[Process unit]{\includegraphics[width=60mm]{images/ch2/triplereptop.pdf}} \\
 	\subfloat[Triple representation]{\includegraphics[width=146mm]{images/ch2/triplerepbot.pdf}} 	
\end{tabular}
\caption{Triple representation of the same energy requirement  \cite{Muller3785THESES}}
\label{fig2:triple}
\vspace{-5mm}
\end{figure*}


\cref{fig2:triple} shows an example of the triple representation for a heat requirement of a process unit. Hot water receives the heat requirement from the steam and delivers it to the process. The condensate receives the same amount of heat from combusted gases in the boiler to become steam. The exergy losses can be visualized in \cref{fig2:triple}b at the right hand side, using the Carnot Composite Curves \cite{LENI-ARTICLE-1996-002}.

\subsection{Site scale integration}
A large industrial process system can be analyzed by either \textit{Single Process Integration} (SPI) or through the \textit{Total Site Integration} (TSI). Both SPI and TSI lay on the same line in terms of approach, but in TSI the opportunity of heat recovery between the process systems is considered rather for the whole site than for each process system separately (SPI). In other words, the process systems in SPI are considered as restricted matches since no heat recovery between the process systems has been considered. We expect therefore that the TSI raises considerably higher potentials of energy saving than the SPI since there is an additional potential of integration between different plants. However, SPI has also its own advantages: it is much easier to implement and would impose less investment cost since it discourages the long-distance connection between different plants. As discussed in \cref{chap2} one of the preliminary challenges faced for total site heat integration is data extraction. A multi-level data extraction has been accordingly introduced in \cref{chap2} based on the different levels of energy requirement analysis to tackle this complexity. A similar idea can also be extended and applied on TSI. This time it would be each process system introduced with a different energy requirement representation, depending on the most detail accessible level of data for that specific process system. Consequently, the TSI could be performed now with a combination of different levels of energy requirement. 



\subsection{Caste study II: Site-Scale integration}

A large industrial process system can be analyzed by either \textit{Single Process Integration} (SPI) or through the \textit{Total Site Integration} (TSI). Both SPI and TSI lay on the same line in terms of approach, but in TSI the opportunity of heat recovery between the process systems is considered rather for the whole site than for each process system separately (SPI). In this case study, the SPI and TSI have been performed for the three process systems (Called as A, B and C) of a site by considering different representations***** for each process system. The total heat requirement of the site is considered as the basis of 100\% unit here, and all the integration analysis values are based on this reference. \cref{fig3:new} shows the repartition of the total consumption between the three process systems. It should be noticed that this 100\% unit does not define the heat consumption of the site cold streams because it considers the present heat recovery in the processes. The major consumers of each process system are identified such that 80\% of the total site consumption are considered in the analysis (Pareto principle). The remaining 20\% is therefore represented by black-box model and will not be considered for heat recovery in the analysis. 

\begin{figure*}[!ht]
\vspace{5mm}
\begin{center}
\includegraphics[width=125mm]{images/ch3/integrationAnalysis.pdf}
\caption{Summary of SPI and TSI for the three process subsystems of the site}
\label{fig3:new}
\end{center}
\end{figure*}

\subparagraph{Single Process Integration (SPI)}

The SPI is applied to each of the three process systems (\cref{fig3:III,fig3:IV,fig3:V}). The SPI analysis for the process system A is performed with a simple-model level***** of detail. 13\% of the total site consumption is explained by process system A, in which, after Pareto analysis, 11\% unit are considered in the integration analysis. The Composite Curves of this process system (\cref{fig3:III}b) show an increase for the potential of heat recovery from 2\% unit in the present situation, up to 8\% unit by integration. This means an energy saving of 62\% for the process system A, but only 7\% for the total site consumption. It should be noted that the 2\% unit heat recovery by P-P***** heat exchangers do not appear in the black-box model*****.

\vspace{10pt}
\begin{figure*}[!ht]
\begin{center}
\includegraphics[width=146mm]{images/ch3/fig2ab.pdf}
\caption{Composite Curves of process system A from black-box to simple-model analysis with SPI (100\% corresponds to the present total site consumption)}
\label{fig3:III}
\end{center}
\end{figure*}

The same strategy is applied to the process system B, which represents 80\% of the whole consumption of the site. Since more detailed data were available from simulation model of this process system, the energy requirement analysis is using the detailed-model ****(\cref{fig3:IV}) that explains 62\% of the heat requirement after the Pareto analysis. The heat recovery potential is increased from 30 with the present heat exchangers, up to 47\% unit after integration. This corresponds to an energy saving of 17\% of the total site consumption or 19\% of process system B consumption.

 
\begin{figure*}[!ht]
\begin{center}
\includegraphics[width=146mm]{images/ch3/fig3ab.pdf}
\caption{Composite Curves of process system B from  black-box to detailed-model analysis with SPI (100\% corresponds to the present total site consumption)}
\label{fig3:IV}
\end{center}
\end{figure*}

Quite the opposite, the available data for process system C, which represents 7\% of the whole site heat consumption, is limited to the energy bill. It is therefore represented by the  black-box model***** (\cref{fig3:V}) and remains at the utility representation level**** ****.\footnote{This is another practical instance of how our multi-level data gathering procedure can come helpful to tune the analysis level of detail with respect to the available data} 
 
\vspace{10pt}
\begin{figure}[!ht]
\begin{center}
\includegraphics[width=80mm]{images/ch3/fig4.pdf}
\caption{Composite Curves of process system C with black-box analysis (100\% corresponds to the present total site consumption)}
\label{fig3:V}
\end{center}
\end{figure}

 Summing up the possible energy savings for the three process systems by SPI shows an additional heat recovery potential of 23\% unit. \cref{fig3:new} gives the detail of SPI analysis of the three process systems.
 
\subparagraph{Total Site Integration (TSI)}

In the TSI approach, combined with the multi-level energy requirement description, the first composite curves (\cref{fig3:VI}a) have been generated by applying the black-box analysis in which the utility representation of all process systems of the site is presented (after Pareto analysis). The result of this analysis gives a general overview of the current energy requirements of the site. However, the heat recovery potential cannot be pointed out with this representation. Consequently, the composite curves are systematically defined by upgrading the T-H profiles of the process systems with a more detailed requirement representation level.
 
 \begin{figure*}[!ht]
\vspace{5mm}
\begin{center}
\includegraphics[width=146mm]{images/ch3/fig5.pdf}
\caption{Systematic improvement of Composite and Grand Composite Curves of total site by combining different energy requirement levels (100\% corresponds to the present total site consumption)}
\label{fig3:VI}
\end{center}
\end{figure*}

 \cref{fig3:VI}c shows the situation where the process system A has been represented by the simple-model while the other process systems remain with the utility representation. Upgrading the process system A from black-box to simple-model analysis, shows a heat recovery potential of 9\% unit that corresponds an additional 1\% unit recovered by exchanging heat between process system A and the other process systems. The analysis is continued by upgrading process system B from black-box to detailed-model (\cref{fig3:VI}e). The result of this integration shows that the total site heat requirement can be reduced to 70\% unit (50+20) which corresponds to 30\% energy saving. Comparing the result of TSI with SPI conclusion (Table. \ref{tab3:one}), exhibits 7\% unit more reduction in total heat requirement. This means that the heat requirement of the process system A (5\% unit) can be fully supplied by heat recovery with process system B. In addition, 2\% unit of process system C requirement can also be satisfied by heat recovery. The summary of SPI and TSI has been presented in Table. \ref{tab3:one}. 
 
 \subparagraph{Heat recovery improvement potentials}

With the analysis of the GCC of the site (\cref{fig3:VI}f) the maximum heat recovery can be further improved by either modifying the process operating conditions (increasing/decreasing the pressure of columns) or by adding supplementary equipment (MVR or HP). The purpose of all categories of modification is to move hot streams from  below to above the pinch or to move the cold streams located above the pinch to below the pinch (plus-minus principle). 

\begin{figure}[!ht]
\centering
\includegraphics[width=71mm]{images/ch3/figpressa.pdf}
\includegraphics[width=70mm]{images/ch3/figpressb.pdf}
\caption{TSI for the process systems with and without pressure modification}
\label{fig3:VII}
\end{figure}

 Since the  streams that created the pinch are the condenser and reboiler of two distillation column units, the first proposition would be to modify the pressure of those columns. Therefore, the pressure of column (1) is modified ($\Delta p=0.4$ bar), in order to increase the temperature of its condenser up to above of the re-boiler of the column (2) (see \cref{fig3:VII}a). Consequently, the heat requirement is reduced by 8\% unit, therefore the overall hot MER is reduced to 42\% unit. Another possibility to upgrade the temperature of the condenser (1) above the re-boiler (2) is to apply heat pumping systems. The HPs can be implemented in two different ways: it can either be an HP with a closed cycle which uses an external fluid as a heat transfer intermediate, or an MVR which uses the process stream as a working fluid of the cycle. 
 
 \begin{figure}[!ht]
\centering
\includegraphics[width=71mm]{images/ch3/figmvra.pdf}
\includegraphics[width=70mm]{images/ch3/figmvrb.pdf}
\caption{TSI for the process systems with and without  MVR and heat pumps}
\label{fig3:VIII}
\end{figure}

 Considering the conditions of pinch streams (stream phases and temperature range), a multi-stage MVR is integrated to the column (1) in order to increase the temperature of the condenser (1) above of the re-boiler (1) and reboiler (2) (see \cref{fig3:VII}b). It should be mentioned that column (1) and (2) are in the same process system. Furthermore, in order to increase the recompression potential and also to avoid a quick creation of new pinch points, two HP systems are integrated to the lower edge of the composite curve. The heat requirement of the site is further reduced by 23\% unit with the integration of multi-stage MVR and HPs, for the expense of 6\% unit mechanical power requirement . The site energy bill is therefore reduced to 47\% unit (27+20). The optimal flow rates in the heat pumping system are obtained by optimization and integration simultaneously with the energy conversion unit. The optimal integration of MVR and HPs is therefore performed together with utility systems. 
 
 \subparagraph{Integration and optimization of energy conversion units}

 \cref{fig3:VIIII} shows the site utility integration and optimization results with Carnot factor representation\footnote{The ambient temperature ($T_0$) is assumed to be 25 $^{\circ}\mathrm{C}$ } for the three following scenarios: (a) base-case, (b) with pressure modification and (c) MVR and heat pumps integrated. The boiler, cogeneration unit with gas turbine, steam network with steam turbines and refrigeration cycle are the selected units within all three scenarios. One of the waste stream from the process is burned as a fuel in the boiler which creates 7\% unit of the current site heat requirement (0.2 relative heat load in \cref{fig3:VIIII}) and corresponds to the highest temperature level of utility streams. A gas turbine produces electricity with an efficiency of 33\% and produces a hot stream from 558 $^{\circ}\mathrm{C}$ to the stack temperature of 160 $^{\circ}\mathrm{C}$ with a thermal efficiency of 50\%. In the calculation, the part load efficiency of the gas turbine is not considered, however, the level of usage is verified to be consistent with its capacity. The price for different forms of energy, influences the choice of the appropriate technology and its level of usage in the integrated solution. Here, the ratio of gas price to the electricity is supposed to be $0.8$. The cogeneration unit is quite sensitive to the price of electricity and gas, thus in the presented solutions, the maximum capacity of the cogeneration system is optimized to produce an equivalent amount heat to the MER of each scenario. As a result, the electricity produced by the cogeneration unit, comes only from the excess exergy content of the necessary required gas that satisfies MER. 
 
  \vspace{5mm}
 \begin{figure}[!ht]
 \begin{center}
 \includegraphics[width=70mm]{images/ch3/figutilitya.pdf}
 \includegraphics[width=70mm]{images/ch3/figutilityb.pdf}
 \includegraphics[width=70mm]{images/ch3/figutilityc.pdf}
 \caption{Site utility integration and optimization for three different scenarios}
 \label{fig3:VIIII}
 \end{center}
  \vspace{-5mm}
 \end{figure}
 
 The proposed solutions are simplified based on both technical feasibilities and the preference of the plant, even if there were a further possibility to produce electricity by the cogeneration unit and through improvement of the steam network. The finalized solutions are then evaluated based on the operating and investment cost and are also compared to the current situation of the plant. The investment cost includes the one required to purchase new equipments, such as the heat exchangers, compressors and turbines. The summary of the analysis is presented in \cref{tab3:two}. The reported relative electricity balance is using the current natural gas consumption as a reference and is positive when electricity is exported. It accounts only for the utility system balance and does not include electricity consumption of the processes. The relative annualized investment cost similarly is using the current annual operating cost as the reference. For the estimation of the investment cost, standard techniques have been used. 
 
 \label{tab3:two}
 
 This analysis shows that the natural gas and electricity consumption have been reduced in each scenario. 
Since the current site uses a simple boiler, installation of a cogeneration system slightly increases the natural gas consumption, simultaneously satisfying MER and producing electricity. However, the total cost of the TSI scenario is reduced considerably by 22\% due to the cogeneration and heat recovery effects. The reduction of gas consumption in the scenario with MVR and heat pumps reaches 40\%, but its effect on the total cost is only 30\% due to the electricity consumption of MVR system and its higher investment cost.
 
 
